\textbf{Duración:} 3 horas (en aula)\\
\textbf{Enfoque:} Deck mínimo viable; notas de orador; variantes clase/conferencia; accesibilidad básica.

\subsubsection*{Objetivos de la sesión}
Al finalizar, la persona participante:
\begin{enumerate}
\item Construye un deck MV (8--12 diapositivas) a partir del storyboard.
\item Genera notas de orador alineadas a la voz docente.
\item Aplica una mini-checklist de accesibilidad (contraste, ritmo, alt text conceptual).
\item Declara el microartefacto final EIC hacia el cierre del CADi.
\end{enumerate}

\subsubsection*{Agenda (resumen)}
\begin{tabularx}{\textwidth}{@{}lclX@{}}
\toprule
Bloque & Tiempo & Contenido \\
\midrule
Recap storyboards & 10 min & Priorizar 8--12 slides \\
Demo deck + notas & 25 min & Voz docente y variantes \\
Accesibilidad & 20 min & Contraste, ritmo, notación legible \\
Taller & 70 min & Deck MV + notas \\
Mini-revisión & 25 min & Checklist accesibilidad \\
Cierre & 20 min & Definir microartefacto final \\
\bottomrule
\end{tabularx}

\subsubsection*{Producto esperado}
Deck MV + notas de orador + mini-checklist de accesibilidad (E6).

\subsubsection*{Trabajo independiente relacionado}
Sesión 7: artículos --- estructura, borrador y citas. TI: avance microartefacto (C) + plantillas (B).
